\section{Introdução}
\label{sec:introducao}

% A poluição atmosférica urbana, em particular a exposição prolongada ao dióxido de nitrogênio (NO\textsubscript{2}) e ao material particulado fino (PM\textsubscript{2.5}), está associada a impactos relevantes sobre mortalidade e morbidade, tornando a estimação precisa desses poluentes um requisito central em avaliações de impacto em saúde pública \cite{who_air_pollution_2022}. Nesse contexto, cresce o interesse por redes de sensores de baixo custo como alternativa complementar às estações de monitoramento de referência, pois permitem maior resolução espacial e temporal, embora apresentem medições sujeitas a deriva, sensibilidade cruzada e forte dependência de parâmetros meteorológicos. Isso exige calibração robusta baseada em modelos estatísticos supervisionados. Evidências empíricas demonstram que a integração de dados ambientais e medições de referência via regressões multivariadas pode melhorar substancialmente a concordância com estações oficiais, destacando a necessidade de ajustes específicos para cada sensor e ambiente operacional \cite{vanzoest2019}.

% A literatura recente também indica uma consolidação de técnicas de aprendizado de máquina aplicadas à análise e previsão de poluentes atmosféricos. Em ambientes internos, modelos recorrentes do tipo RNN–LSTM foram empregados para prever concentrações de NO\textsubscript{2} em dormitórios estudantis, demonstrando a capacidade dessas arquiteturas em capturar dependências temporais de médio e longo prazo \cite{xie2025}. Em escala de sensores miniaturizados, modelos híbridos combinando análise stepwise e perceptron multicamada têm se mostrado superiores a modelos puramente lineares em tarefas de calibração, ao modelarem adequadamente relações não lineares entre preditores e respostas \cite{liu2021}. Esses resultados reforçam o potencial do uso combinado de técnicas estatísticas clássicas e aprendizado profundo na avaliação da qualidade do ar.

% No tocante à modelagem de poluentes atmosféricos a partir de parâmetros meteorológicos, estudos comparativos entre Regressão Linear Múltipla, técnicas baseadas em componentes (PCR, PLS) e regressões penalizadas (Ridge) indicam que estas últimas tendem a apresentar melhor estabilidade numérica em cenários com forte multicolinearidade, mantendo elevada capacidade preditiva \cite{polat2021}. Adicionalmente, investigações recentes mostram que diferentes configurações de modelos lineares podem maximizar o equilíbrio entre ajuste, robustez e interpretabilidade em análises de potencial oxidativo de partículas inaláveis \cite{ngocthuy2024}.

% Este trabalho insere-se nesse cenário ao investigar a calibração multivariada de sensores de baixo custo para NO\textsubscript{2} utilizando o conjunto de dados Air Quality UCI \cite{Souza2025EDA,devito2008}. O dataset inclui respostas de sensores MOx, variáveis meteorológicas e medidas de referência de diversos poluentes em ambiente urbano monitorado. Definimos a concentração de NO\textsubscript{2} real (NO2(GT)) como variável alvo e utilizamos as leituras dos sensores e fatores meteorológicos como preditores, reproduzindo um cenário típico de redes urbanas de monitoramento com sensores de baixo custo.

% Metodologicamente, executamos um pipeline comparativo com pré-processamento rigoroso, incluindo remoção de valores ausentes, transformação logarítmica para variáveis assimétricas e padronização Z-score com parâmetros estimados apenas no conjunto de treino. Foram implementados e avaliados modelos OLS, Ridge e PLS, com seleção de hiperparâmetros via validação cruzada K-fold, além de um MLP como referência não linear, visando mensurar os ganhos efetivos decorrentes do aumento de complexidade.

% Os resultados obtidos indicam que modelos lineares adequadamente regularizados apresentam desempenho competitivo, com erros médios quadráticos na faixa de 10–11 unidades de NO\textsubscript{2} e coeficientes de determinação próximos de 0,88–0,90. PLS e Ridge apresentaram desempenho mais estável que OLS, enquanto o MLP não superou consistentemente as alternativas lineares. Também observamos que canais associados a NOx e sensores PT08 exibem maior contribuição para a calibração. Esses achados sugerem que uma combinação de pré-processamento cuidadoso, regressões penalizadas e componentes supervisionados oferece um compromisso favorável entre acurácia, interpretabilidade e custo computacional, alinhado às recomendações recentes da literatura especializada \cite{liu2021,xie2025,vanzoest2019,polat2021,ngocthuy2024}.
